% ======================================================================
% Chapter: The Level-1 KM/Lattice Bridge
% ======================================================================

\chapter{The level-1 bridge}\label{chap:level1-bridge}

At level~$1$, two independently constructed bar complexes produce
identical cohomology by the Frenkel--Kac--Segal theorem
(Theorem~\ref{thm:lattice:frenkel-kac}). The affine Kac--Moody algebra
$\widehat{\fg}_1$ at level~$1$ and the lattice vertex algebra
$\Vlat_{\Lambda_\fg}$ on the root lattice of a simply-laced~$\fg$
are the same algebra in two presentations: the current presentation
(with $\dim\fg$ generating fields $J^a(z)$) and the lattice
presentation (with $\operatorname{rank}(\fg)$ Heisenberg fields
and vertex operators $E_\alpha(z) = Y(e^\alpha, z)$).
The bar constructions built from these two descriptions must
agree, and the comparison reveals a structural phenomenon:
the modular characteristic $\kappa$ undergoes a
\emph{level-$1$ reduction} from the generic-level Kac--Moody
value to the lattice rank.

\begin{table}[ht]
\centering
\small
\caption{Five-theorem verification for $\widehat{\fg}_1
\cong \Vlat_{\Lambda_\fg}$ (simply-laced).}%
\label{tab:level1-five-theorems}
\begin{tabular}{lp{5.8cm}ll}
\toprule
\textbf{Theorem} & \textbf{Statement}
 & \textbf{Status} & \textbf{Reference} \\
\midrule
A (adjunction) &
 $(\widehat{\fg}_1)^!
 = \mathrm{CE}^{\mathrm{ch}}(\widehat{\fg}_{-1-2h^\vee})$;
 equiv.\ $(\Vlat_{\Lambda_\fg}^{\varepsilon})^!
 = (\Vlat_{\Lambda_\fg}^{\varepsilon^{-1}})^c$
 & Proved
 & Thms~\ref{thm:universal-kac-moody-koszul},
 \ref{thm:lattice:koszul-dual} \\
B (inversion) &
 $\Omega(\barB(\widehat{\fg}_1))
 \xrightarrow{\sim} \widehat{\fg}_1$
 & Proved
 & Thm~\ref{thm:sl2-genus1-inversion} \\
C (complementarity) &
 $Q_g(\widehat{\fg}_1) \oplus
 Q_g(\widehat{\fg}_{-1-2h^\vee}) = H^*(\overline{\cM}_g, Z)$
 & Proved
 & Thm~\ref{thm:sl2-genus1-complementarity} \\
D (modular char.) &
 $\kappa(\widehat{\fg}_1) = \operatorname{rank}(\fg)$
 (Prop.~\ref{prop:level1-kappa-reduction})
 & Proved
 & \S\ref{sec:level1-kappa} \\
H (Hochschild) &
 ChirHoch${}^*(\widehat{\fg}_1)$ polynomial
 & Proved
 & Thm~\ref{thm:kac-moody-ainfty} \\
\bottomrule
\end{tabular}
\end{table}

\begin{table}[ht]
\centering
\small
\caption{Shadow archetype data for $\widehat{\fg}_1
\cong \Vlat_{\Lambda_\fg}$.}%
\label{tab:level1-shadow-archetype}
\begin{tabular}{ll}
\toprule
\textbf{Invariant} & \textbf{Value} \\
\midrule
Class & G (Gaussian) --- reduced from L at $k = 1$ \\
Shadow depth $r_{\max}$ & $2$ \\
$\kappa$ & $\operatorname{rank}(\fg)$ \\
Cubic shadow $\mathfrak{C}$ & $0$ (Jacobi + null-vector exactness) \\
$r$-matrix $r(z)$ &
 $\Omega/\bigl((1{+}h^\vee)\,z\bigr)$ (KZ normalisation at $k = 1$; cf.\ Remark~\ref{rem:km-collision-residue-rmatrix}) \\
Koszul dual &
 $\mathrm{CE}^{\mathrm{ch}}(\widehat{\fg}_{-1-2h^\vee})
 \;=\; (\Vlat_{\Lambda_\fg}^{\varepsilon^{-1}})^c$ \\
Complementarity &
 $\kappa + \kappa' = 0$ (affine KM) \\
$L_\infty$-formality &
 Strictly formal (depth $2$, all shadows vanish for $r \geq 3$) \\
\bottomrule
\end{tabular}
\end{table}


% ======================================================================
% Section 1: The FKS isomorphism
% ======================================================================

\section{The Frenkel--Kac--Segal isomorphism at the chain level}%
\label{sec:level1-fks}

The Frenkel--Kac--Segal theorem
(Theorem~\ref{thm:lattice:frenkel-kac}; \cite{FK80,Se81})
establishes a vertex algebra isomorphism
\begin{equation}\label{eq:level1-fks}
\varphi_{\mathrm{FKS}} \colon
L(\widehat{\fg}, 1)
\;\xrightarrow{\;\cong\;}
\Vlat_{\Lambda_\fg}^{\varepsilon}
\end{equation}
for every simply-laced simple Lie algebra~$\fg$. At level~$1$
simply-laced, the universal vacuum module~$V_1(\fg)$ is simple
(no null vectors in the vacuum Verma module), so
$V_1(\fg) = L(\widehat{\fg}, 1) = \Vlat_{\Lambda_\fg}$.

The isomorphism $\varphi_{\mathrm{FKS}}$ sends
each Cartan current $H_i(z)$ to the Heisenberg field
$Y(\alpha_{i,-1}|0\rangle, z)$
and each root current $E_\alpha(z)$ to the vertex
operator $Y(e^\alpha, z)$. In particular, the
$\dim\fg$ generators of the current presentation are not
independent at the lattice level: the $\dim\fg - \operatorname{rank}(\fg)$
root currents are vertex operators of Heisenberg
composites~$e^\alpha$.

\begin{remark}[Two generating sets, one algebra]%
\label{rem:level1-two-presentations}
The current algebra~$\widehat{\fg}_1$ is freely strongly
generated by $\dim\fg$ fields of weight~$1$: the Cartan
currents $\{H_i\}_{i=1}^r$ and the root
currents $\{E_\alpha\}_{\alpha \in \Phi}$.
The lattice VOA~$\Vlat_{\Lambda_\fg}$ is freely strongly
generated by $r = \operatorname{rank}(\fg)$ Heisenberg fields
$\{h_i\}_{i=1}^r$, with the vertex operators $E_\alpha(z)$
arising as normally ordered exponentials of the $h_i$.
The bar complex is an invariant of the algebra, not of a
presentation: the two generating sets produce
quasi-isomorphic bar complexes. The comparison of the
resulting bar-level invariants is the content of this chapter.
\end{remark}


% ======================================================================
% Section 2: The kappa reduction
% ======================================================================

\section{The modular characteristic at level~$1$}%
\label{sec:level1-kappa}

At generic level~$k \neq -h^\vee$, the modular characteristic
of the affine Kac--Moody algebra is
(equation~\eqref{eq:km-kappa-derivation}):
\begin{equation}\label{eq:level1-km-kappa-generic}
\kappa(\widehat{\fg}_k)
= \frac{\dim(\fg)\,(k + h^\vee)}{2h^\vee}.
\end{equation}
The lattice curvature-braiding orthogonality theorem
(Theorem~\ref{thm:lattice:curvature-braiding-orthogonal})
gives:
\begin{equation}\label{eq:level1-lattice-kappa}
\kappa(\Vlat_{\Lambda_\fg})
= \operatorname{rank}(\Lambda_\fg)
= \operatorname{rank}(\fg).
\end{equation}
At $k = 1$ simply-laced, both formulas apply to the
same algebra. Their values are:
\begin{center}
\renewcommand{\arraystretch}{1.3}
\begin{tabular}{lccc}
\toprule
$\fg$ & $\kappa^{\mathrm{KM}}
 = \frac{\dim\fg\,(1+h^\vee)}{2h^\vee}$
 & $\kappa^{\mathrm{lat}} = \operatorname{rank}(\fg)$
 & $c = \frac{\dim\fg}{1+h^\vee}$ \\
\midrule
$\mathfrak{sl}_2$
 & $9/4$
 & $1$
 & $1$ \\
$\mathfrak{sl}_3$
 & $8 \cdot 4/6 = 16/3$
 & $2$
 & $2$ \\
$D_4$
 & $28 \cdot 7/12 = 49/3$
 & $4$
 & $4$ \\
$E_6$
 & $78 \cdot 13/24 = 169/4$
 & $6$
 & $6$ \\
$E_7$
 & $133 \cdot 19/36 = 2527/36$
 & $7$
 & $7$ \\
$E_8$
 & $248 \cdot 31/60 = 1922/15$
 & $8$
 & $8$ \\
\bottomrule
\end{tabular}
\end{center}
The two columns disagree. The lattice value
$\kappa^{\mathrm{lat}} = \operatorname{rank}(\fg)$ is
correct; the Kac--Moody formula~\eqref{eq:level1-km-kappa-generic}
overcounts because its derivation treats all
$\dim\fg$ currents as \emph{independent} generators
of the bar complex.

\begin{proposition}[Level-$1$ $\kappa$ reduction; \ClaimStatusProvedHere]%
\label{prop:level1-kappa-reduction}%
\index{modular characteristic!level-1 reduction}%
\index{kappa@$\kappa$!level-1 bridge}
For $\fg$ simply-laced simple and $k = 1$:
\begin{equation}\label{eq:level1-kappa-reduction}
\kappa(\widehat{\fg}_1)
= \operatorname{rank}(\fg),
\end{equation}
and not $\dim(\fg)(1+h^\vee)/(2h^\vee)$. The generic-level
formula~\eqref{eq:level1-km-kappa-generic} applies to the
universal vacuum module~$V_k(\fg)$ with $\dim\fg$
independent weight-$1$ generators; at $k = 1$, the lattice
realization identifies the root currents as composites of the
$\operatorname{rank}(\fg)$ Heisenberg generators, reducing
the effective generator count and thereby the modular
characteristic.
\end{proposition}

\begin{proof}
By the Frenkel--Kac--Segal isomorphism
$\varphi_{\mathrm{FKS}} \colon V_1(\fg)
\xrightarrow{\cong} \Vlat_{\Lambda_\fg}$,
the modular characteristic $\kappa$ is computed from either
presentation. The lattice computation
(Theorem~\ref{thm:lattice:curvature-braiding-orthogonal})
decomposes the genus-$1$ bar complex into Cartan and root
sectors. The Cartan sector (zero lattice charge) has
$\dfib^{\,2} = \operatorname{rank}(\fg) \cdot \omega_1$,
contributed by $\operatorname{rank}(\fg)$ independent
Heisenberg bosons. Each root sector (lattice charge
$\gamma = \alpha_i + \alpha_j$ with
$\langle \alpha_i, \alpha_j \rangle = -1$) has
$\dfib^{\,2} = 0$: the simple-pole residue is unchanged at
genus~$1$
(Proposition~\ref{prop:lattice:genus1-simple-pole}), and
the root-sector bar differential is identical to its
genus-$0$ value. The total curvature is the Cartan
contribution alone: $\kappa = \operatorname{rank}(\fg)$.

The discrepancy with the
generic-level formula arises because
$\kappa^{\mathrm{KM}} = \dim(\fg)(k+h^\vee)/(2h^\vee)$
is derived from the OPE of $\dim\fg$ currents treated as
\emph{independent} generators
(equation~\eqref{eq:km-dfib-squared-trace}). This derivation
includes a simple-pole contribution
$\kappa_{\mathrm{sp}} = \dim(\fg)/2$ from the structure
constants $f^{abc}$. At generic level, this contribution is
a genuine part of the bar curvature because the root currents
are algebraically independent of the Cartan currents. At
$k = 1$ simply-laced, the FKS isomorphism identifies each root
current $E_\alpha(z) = Y(e^\alpha, z)$ as a vertex operator
built from the Heisenberg fields. These vertex operators are
composites, not independent generators: the bar complex of
$\Vlat_{\Lambda_\fg}$ factors through the lattice
weight filtration, and the root-sector curvature
$\dfib^{\,2}|_{\gamma \neq 0} = 0$ by the curvature-braiding
orthogonality theorem. The simple-pole contribution
$\kappa_{\mathrm{sp}}$ that appears in the universal-level
derivation is killed at $k = 1$ by the lattice structure.
\end{proof}

\begin{remark}[The central charge coincidence]%
\label{rem:level1-central-charge}%
At level~$1$ simply-laced, the central charge satisfies
$c = k\,\dim\fg/(k+h^\vee) = \dim\fg/(1+h^\vee) =
\operatorname{rank}(\fg)$. Thus
$\kappa = c = \operatorname{rank}(\fg)$ at $k = 1$.
This is a coincidence specific to level~$1$: at generic~$k$,
$\kappa \neq c$ and $\kappa \neq c/2$ for affine Kac--Moody
algebras (the discriminant $(k+h^\vee)^2 - kh^\vee > 0$
has no real solutions; see equation~\eqref{eq:km-kappa-derivation}
step~4(x)). The identity $\kappa = c$ at $k = 1$ is an
independent confirmation: it follows from
$\kappa = \operatorname{rank}(\fg)$ (this proposition) and
$c = \operatorname{rank}(\fg)$ (the Sugawara formula at
$k = 1$).
\end{remark}


% ======================================================================
% Section 3: Shadow class reduction
% ======================================================================

\section{Shadow class reduction: from L to G}%
\label{sec:level1-shadow}

At generic level, the affine Kac--Moody algebra
$\widehat{\fg}_k$ is class~L (Lie/tree, $r_{\max} = 3$):
the Lie bracket $[J^a, J^b] = f^{ab}_c J^c$ produces a
nonzero cubic shadow~$\mathfrak{C}$, and the Jacobi identity
kills the quartic obstruction $o_4 = 0$
(Table~\ref{tab:km-shadow-archetype}). At $k = 1$
simply-laced, a stronger vanishing holds.

\begin{proposition}[Cubic shadow vanishing at level~$1$;
\ClaimStatusProvedHere]%
\label{prop:level1-cubic-vanishing}%
\index{shadow obstruction tower!level-1 vanishing}%
\index{cubic shadow!vanishing at level 1}
For $\fg$ simply-laced and $k = 1$:
\begin{equation}\label{eq:level1-cubic-vanishing}
\mathfrak{C}(\widehat{\fg}_1) = 0.
\end{equation}
The shadow obstruction tower terminates at degree~$2$, and
$\widehat{\fg}_1 \cong \Vlat_{\Lambda_\fg}$ is
class~$\mathbf{G}$ \textup{(}Gaussian\textup{)}.
\end{proposition}

\begin{proof}
By Theorem~\ref{thm:lattice:curvature-braiding-orthogonal}(ii),
every root sector of the genus-$1$ bar complex has
$\dfib^{\,2} = 0$. The cubic shadow is the degree-$3$
component of the MC element $\Theta_A$ in the modular
convolution algebra; it receives contributions from the
bar differential acting on $3$-fold tensor products of
generators. In the lattice presentation, the
generating fields are the $\operatorname{rank}(\fg)$ Heisenberg
currents $h_i(z)$, which have only double-pole OPEs
$h_i(z)\,h_j(w) \sim A_{ij}/(z-w)^2$. Since the cubic
shadow requires a simple pole (the Lie bracket), and
Heisenberg generators have no simple poles, the
degree-$3$ component vanishes identically:
$\mathfrak{C}(\Vlat_{\Lambda_\fg}) = 0$.

In the current presentation, the same vanishing holds by
the FKS isomorphism: the cubic shadow of $\widehat{\fg}_1$
computed from the $J^a$ currents must agree with the cubic
shadow of $\Vlat_{\Lambda_\fg}$ computed from the $h_i$
fields (since $\kappa$ and the full shadow tower are
intrinsic to the algebra). The mechanism is that the
level-$1$ null vectors force the Lie bracket contributions
to become exact in the bar complex: the structure constant
$f^{abc}$ contributions to the degree-$3$ MC equation, which
are nonzero at generic~$k$, degenerate into
coboundaries at $k = 1$ via the lattice relations
$E_\alpha = Y(e^\alpha, z)$.
\end{proof}

\begin{remark}[Classification transition at $k = 1$]%
\label{rem:level1-class-transition}%
\index{shadow class!transition at level 1}
The shadow class of $\widehat{\fg}_k$ undergoes a transition
at $k = 1$:
\begin{center}
\begin{tabular}{lcc}
\toprule
\textbf{Level} & \textbf{Class} & $r_{\max}$ \\
\midrule
Generic $k \neq -h^\vee$ & L (Lie/tree) & $3$ \\
$k = 1$ (simply-laced) & G (Gaussian) & $2$ \\
$k = -h^\vee$ (critical) & uncurved ($\kappa = 0$) & $2$ \\
\bottomrule
\end{tabular}
\end{center}
At generic level, the Lie bracket is the source of the cubic
shadow. At the two exceptional levels, the cubic shadow
vanishes for different reasons: at $k = -h^\vee$ the entire
curvature vanishes ($\kappa = 0$);
at $k = 1$ simply-laced, the lattice realization makes the
Lie bracket contributions exact.
\end{remark}


% ======================================================================
% Section 4: Explicit verification
% ======================================================================

\section{Explicit verification: \texorpdfstring{$\mathfrak{sl}_2$}{sl2}
at level~$1$}%
\label{sec:level1-sl2}

For $\fg = \mathfrak{sl}_2$: $\dim\fg = 3$, $h^\vee = 2$,
$\operatorname{rank} = 1$, root lattice $\Lambda = A_1 = \Z\alpha$ with
$\langle\alpha,\alpha\rangle = 2$.

\emph{Central charge.}
$c = 3 \cdot 1/(1+2) = 1 = \operatorname{rank}(A_1)$.

\emph{Modular characteristic.}
The generic formula gives
$\kappa^{\mathrm{KM}} = 3 \cdot 3/(2 \cdot 2) = 9/4$. The lattice
formula gives $\kappa^{\mathrm{lat}} = 1$. By
Proposition~\ref{prop:level1-kappa-reduction},
$\kappa(\widehat{\mathfrak{sl}}_{2,1}) = 1$.

\emph{Genus-$1$ free energy.}
$F_1 = \kappa/24 = 1/24$, the free energy of a single free boson.

\emph{Bar complex.}
In the lattice presentation, $\Vlat_{A_1}$ has a single Heisenberg
generator $h(z)$ of weight~$1$. The bar complex is the Heisenberg bar
complex $\barB(\mathcal{H}_1)$ with $\kappa = 1$. The genus-$1$
curvature is $\dfib^{\,2} = \omega_1$ from the single Cartan boson.
No root-sector curvature arises because
$\langle n\alpha, m\alpha \rangle = -1$ has no integer solutions
(Computation~\ref{comp:lattice:bar-A1}): the lattice bar differential
vanishes on pure lattice-vector elements.

In the current presentation, the three generators
$e(z), h(z), f(z)$ produce a bar complex with
$3$ generators in bar degree~$1$, but the FKS relations
$e = Y(e^\alpha, z)$, $f = Y(e^{-\alpha}, z)$ identify
two of these with vertex operators of the third, and the
resulting quasi-isomorphism reduces the effective bar
complex to the rank-$1$ Heisenberg model.


% ======================================================================
% Section 5: The ADE family
% ======================================================================

\section{The ADE family}%
\label{sec:level1-ade}

\begin{computation}[Level-$1$ bridge data for the simply-laced
series; \ClaimStatusProvedHere]%
\label{comp:level1-ade-bridge}%
\index{level-1 bridge!ADE data}
\begin{center}
\renewcommand{\arraystretch}{1.3}
\begin{tabular}{lccccc}
\toprule
$\fg$ & $\operatorname{rank}$ & $\dim\fg$
 & $c = \operatorname{rank}$
 & $\kappa = \operatorname{rank}$
 & $|\Phi|$ \\
\midrule
$A_1$ & $1$ & $3$ & $1$ & $1$ & $2$ \\
$A_2$ & $2$ & $8$ & $2$ & $2$ & $6$ \\
$A_3$ & $3$ & $15$ & $3$ & $3$ & $12$ \\
$D_4$ & $4$ & $28$ & $4$ & $4$ & $24$ \\
$D_5$ & $5$ & $45$ & $5$ & $5$ & $40$ \\
$E_6$ & $6$ & $78$ & $6$ & $6$ & $72$ \\
$E_7$ & $7$ & $133$ & $7$ & $7$ & $126$ \\
$E_8$ & $8$ & $248$ & $8$ & $8$ & $240$ \\
\bottomrule
\end{tabular}
\end{center}
In every case: $\kappa = c = \operatorname{rank}(\fg) =
\operatorname{rank}(\Lambda_\fg)$, agreeing with the
lattice formula. The number of roots $|\Phi| = \dim\fg -
\operatorname{rank}(\fg)$ counts the root currents that
become composites under FKS; these contribute zero to the
genus-$1$ bar curvature by
Theorem~\ref{thm:lattice:curvature-braiding-orthogonal}(ii).
\end{computation}

\begin{remark}[Non-simply-laced types]%
\label{rem:level1-non-simply-laced}%
\index{level-1 bridge!non-simply-laced}
For non-simply-laced $\fg$ (types $B_n$, $C_n$, $F_4$, $G_2$),
the Frenkel--Kac--Segal theorem does not apply directly: the
root lattice contains roots of different lengths with
$\langle\alpha,\alpha\rangle \neq 2$
(Remark~\ref{rem:lattice:non-simply-laced}). The level-$1$
algebra $\widehat{\fg}_1$ is not a lattice VOA.
At level~$1$ for non-simply-laced types,
$\kappa = \dim(\fg)(1+h^\vee)/(2h^\vee)$ is the full
generic-level formula (no reduction occurs), and
$\widehat{\fg}_1$ remains class~L. The class reduction
$\mathrm{L} \to \mathrm{G}$ at $k = 1$ is a specifically
simply-laced phenomenon, controlled by the lattice
realization.
\end{remark}

\begin{remark}[Unimodularity and self-duality at $k = 1$]%
\label{rem:level1-unimodular}%
\index{Koszul self-duality!level-1 bridge}
Among the ADE root lattices, only $E_8$ is unimodular
($\Lambda_{E_8} = \Lambda_{E_8}^*$). By
Proposition~\ref{prop:lattice:self-dual-criterion},
$\Vlat_{E_8}$ is the unique Koszul self-dual lattice VOA
in the ADE series:
$(\Vlat_{E_8})^! \cong \Vlat_{E_8}$. This self-duality
is visible in the complementarity data:
$\kappa(\widehat{E}_{8,1}) = 8$,
$\kappa' = \kappa(\widehat{E}_{8,-61}) = -8$
(by the Feigin--Frenkel involution $k \mapsto -1-2h^\vee = -61$),
so $\kappa + \kappa' = 0$ (as required for affine KM by
Proposition~\ref{prop:kappa-anti-symmetry-ff}).
For non-unimodular root lattices ($A_n$, $D_n$, $E_6$, $E_7$),
the Koszul dual involves the inverse cocycle
(Theorem~\ref{thm:lattice:koszul-dual}), and the FKS
bridge identifies this with the Feigin--Frenkel dual
$\widehat{\fg}_{-1-2h^\vee}$
(Theorem~\ref{thm:universal-kac-moody-koszul}).
\end{remark}
